% GENERATED FILE. Edit canonical lessons, then run npm run generate:books.
% canonical-source-hash: fnv1a64:5fb445c9e2d6cd36
% canonical-lessons: RU-C261-chital, RU-C261-prochital, RU-C261-poka, RU-C261-review

\chapter{The Past, Seen Two Ways}
\label{ch:ru-the-past-seen-two-ways}
\hlchaptermodality{261}

\begin{chapteropening}
\textbf{By the end of this chapter:} \emph{I can say what happened and what was going on in the past, choosing between \ru{читал} and \ru{прочитал}, and join two events with \ru{когда} and \ru{пока}.}

Complete RU-C261-review: a past evening told aloud, one scene with \ru{пока} and one finished event with \ru{когда}, each verb in the right aspect.
\end{chapteropening}

\section[\texorpdfstring{\ru{читал}, \ru{читала} … (chitál, chitála …)}{читал, читала … (chitál, chitála …)}]{\ru{читал}, \ru{читала}, \ru{читали} — the past, and who did it}

\label{lesson:RU-C261-chital}



\begin{quote}
Say \textquotedblleft{}recently\textquotedblright{} and \textquotedblleft{}yesterday\textquotedblright{}. (\emph{Nedávno. Vcherá.}) The first came with one past sentence: \emph{Nedávno ya chitál knígu.} This lesson takes that \emph{-\ru{л}} apart.
\end{quote}

\begin{grammarlens}[title={Grammar Lens: -\ru{ть} becomes -\ru{л}}]
Drop the infinitive's \textbf{-\ru{ть}} and put \textbf{-\ru{л}} in its place: \emph{\ru{читать}}, \textbf{\ru{читал}}; \emph{\ru{работать}}, \textbf{\ru{работал}}. The past has no endings for \emph{I} or \emph{you}. It changes for \textbf{who} did it: \textbf{-\ru{л}} for a man, \textbf{-\ru{ла}} for a woman, \textbf{-\ru{ло}} for a neuter noun.

\begin{quote}
\textbf{\ru{Вчера} \ru{я} \ru{читал}.} — \emph{vcherá ya chitál} — Yesterday I read; a man says it. \textbf{\ru{Вчера} \ru{я} \ru{читала}.} — \emph{vcherá ya chitála} — the same, from a woman. \textbf{\ru{Мама} \ru{работала}, \ru{а} \ru{папа} \ru{читал}.} — \emph{máma rabótala, a pápa chitál} — Mum was working and Dad was reading.
\end{quote}

\emph{\ru{Папа}} ends in \emph{-\ru{а}} and still takes \textbf{-\ru{л}}: a father is a man.
\end{grammarlens}

\begin{grammarlens}[title={Grammar Lens: -\ru{ли} for more than one}]
For more than one, gender drops away: \textbf{-\ru{ли}}. \textbf{\ru{Мы} \ru{гуляли} \ru{в} \ru{парке}.} — \emph{my gulyáli v párke} — We went for a walk in the park. The polite \textbf{\ru{вы}} takes \textbf{-\ru{ли}} even for one person: \textbf{\ru{Вы} \ru{читали}?} — \emph{vy chitáli?}
\end{grammarlens}

\begin{grammarlens}[title={Grammar Lens: \ru{был}, \ru{была}, \ru{было}, \ru{были}}]
\emph{\ru{Быть}} is silent in the present, but its past follows the same pattern: \textbf{\ru{был}} (\emph{byl}), \textbf{\ru{была}} (\emph{bylá}), \textbf{\ru{было}} (\emph{býlo}), \textbf{\ru{были}} (\emph{býli}). The stress moves to the end for a woman.

\begin{quote}
\textbf{\ru{Я} \ru{была} \ru{в} \ru{музее}.} — \emph{ya bylá v muzéye} — I was at the museum.
\end{quote}
\end{grammarlens}

\begin{cousinweb}
By the usual historical account, the \emph{-\ru{л}} form was once a describing word, \textquotedblleft{}having read\textquotedblright{}, said beside a present form of \textbf{\ru{быть}}: Old Russian \emph{čitalŭ jesmĭ}. The helping verb wore away and the describing word stayed, so the past still agrees like an adjective. \textbf{\ru{Быть}} goes back to Proto-Indo-European \emph{*b\textsuperscript{h}uH-}, \textquotedblleft{}to become\textquotedblright{} (EtymDB 2.1); English \emph{be} comes from the same root.
\end{cousinweb}

\subsection*{Your turn}
\begin{itemize}
\raggedright
  \item \emph{Say it:} \emph{Vcherá ya chitál}, then \emph{Vcherá ya chitála}
  \item \emph{Say it:} \emph{Máma rabótala, a pápa chitál.}
  \item \emph{Say it:} \emph{My gulyáli v párke.}
  \item \emph{Say it:} \emph{byl, bylá, býlo, býli}
  \item \emph{Answer:} \emph{Gde ty byl vcherá?} \emph{Ya byl} or \emph{Ya bylá}, and a place
\end{itemize}

\begin{tcolorbox}[breakable,colback=teal!4,colframe=teal!35!black,title={Before you move on}]
What replaces \emph{-\ru{ть}} in the past? (\textbf{-\ru{л}.}) What does a woman add? (\textbf{-\ru{а}:} \emph{\textbf{\ru{читала}}}.) More than one person, or polite \emph{\ru{вы}}? (\textbf{-\ru{ли}.}) \textquotedblleft{}She was at the museum\textquotedblright{}? (\emph{\textbf{\ru{Она} \ru{была} \ru{в} \ru{музее}.}})
\end{tcolorbox}
\section[\texorpdfstring{\ru{читал} / \ru{прочитал} (chitál / prochitál)}{читал / прочитал (chitál / prochitál)}]{\ru{читал} or \ru{прочитал} — the reading, and the book read through}

\label{lesson:RU-C261-prochital}



\begin{quote}
Name the partners you met with the first verbs: \emph{\ru{читать}, \ru{прочитать}}; \emph{\ru{писать}, \ru{написать}}. One is the doing, the other the deed done. In the present only the doing is used. In the past both are, and the choice decides the story.
\end{quote}

\begin{grammarlens}[title={Grammar Lens: the doing, or the deed done}]
\begin{quote}
\textbf{\ru{Вчера} \ru{я} \ru{читал} \ru{книгу}.} — \emph{vcherá ya chitál knígu} — Yesterday I was reading a book. \textbf{\ru{Вчера} \ru{я} \ru{прочитал} \ru{книгу}.} — \emph{vcherá ya prochitál knígu} — Yesterday I read the book, to the end.
\end{quote}

The \textbf{imperfective} shows the reading going on, and says nothing about the last page. The \textbf{perfective} sees the deed whole: finished, with a result. The book is read.
\end{grammarlens}

\begin{grammarlens}[title={Grammar Lens: a habit, and a bare fact}]
The imperfective has two more jobs. A \textbf{habit}: \textbf{\ru{Я} \ru{часто} \ru{читал} \ru{газету}.} — \emph{ya chásto chitál gazétu} — I often read the paper. And the \textbf{bare fact} that something happened at all:

\begin{quote}
\textbf{— \ru{Ты} \ru{читал} \ru{эту} \ru{книгу}? — \ru{Да}, \ru{читал}.} — \emph{ty chitál étu knígu? da, chitál} — Have you read this book? Yes, I have.
\end{quote}
\end{grammarlens}

\begin{culture}
\begin{quote}
\textbf{\ru{Мы} \ru{посмотрели} \ru{фильм}.} — \emph{my posmotréli film} — We watched the film, all of it. \textbf{\ru{Я} \ru{писал} \ru{сообщение}, \ru{но} \ru{не} \ru{написал}.} — \emph{ya pisál soobshchéniye, no ne napisál} — I was writing a message, but did not get it written.
\end{quote}
\end{culture}

\begin{cousinweb}
Russian grammar calls aspect \textbf{\ru{вид}}, \textquotedblleft{}a look, a view\textquotedblright{}, which EtymDB 2.1 takes through Proto-Slavic \emph{*vidŭ} to PIE \emph{*weyd-}, \textquotedblleft{}to see, to know\textquotedblright{}: the root of \textbf{\ru{видеть}} and of English \emph{wit}. Two aspects, two views of one deed: from inside, while it goes on, or from outside, whole. The prefix \emph{\ru{про}-} means \textquotedblleft{}through\textquotedblright{}: \textbf{\ru{прочитать}}, to read through.
\end{cousinweb}

\subsection*{Your turn}
\begin{itemize}
\raggedright
  \item \emph{Say it:} \emph{Vcherá ya chitál knígu.} The reading went on
  \item \emph{Say it:} \emph{Vcherá ya prochitál knígu.} The book is finished
  \item \emph{Say it:} \emph{Ya chásto chitál gazétu.} A habit
  \item \emph{Answer:} \emph{Ty chitál étu knígu?} \emph{Da, chitál} or \emph{Net, ne chitál}
  \item \emph{Say it:} \emph{My posmotréli film.}
\end{itemize}

\begin{tcolorbox}[breakable,colback=teal!4,colframe=teal!35!black,title={Before you move on}]
\textquotedblleft{}I finished the book\textquotedblright{}: which verb? (\emph{\textbf{\ru{Прочитал}.}}) \textquotedblleft{}I often read the paper\textquotedblright{}? (\emph{\textbf{\ru{Читал}}}, a habit.) What does Russian call aspect? (\emph{\textbf{\ru{Вид}}}, a view.)
\end{tcolorbox}
\section[\texorpdfstring{\ru{когда} …, …; \ru{пока} … (kogdá …, …; poká …)}{когда …, …; пока … (kogdá …, …; poká …)}]{\ru{когда} and \ru{пока} — the scene, and what cut into it}

\label{lesson:RU-C261-poka}



\begin{quote}
Say \textquotedblleft{}yesterday I was reading a book\textquotedblright{}, then \textquotedblleft{}yesterday I read the book to the end\textquotedblright{}. (\emph{Vcherá ya chitál knígu. Vcherá ya prochitál knígu.})

And the hinge from the present: \emph{kogdá ya chitáyu, ya dúmayu}. \textbf{\ru{Когда}} joins two halves with a comma at the seam. In the past, the aspect on each side tells the listener how the two halves fit together.
\end{quote}

\begin{grammarlens}[title={Grammar Lens: one going on, one cutting in}]
\begin{quote}
\textbf{\ru{Когда} \ru{я} \ru{читал}, \ru{позвонил} \ru{брат}.} — \emph{kogdá ya chitál, pozvoníl brat} — When I was reading, my brother rang.
\end{quote}

The reading is imperfective: it was going on. \textbf{\ru{Позвонил}}, the perfective of \emph{\ru{звонить}}, is one call that cut into it.

Put a perfective on \textbf{both} sides and the events come one after the other:

\begin{quote}
\textbf{\ru{Когда} \ru{я} \ru{прочитал} \ru{книгу}, \ru{я} \ru{позвонил} \ru{брату}.} — \emph{kogdá ya prochitál knígu, ya pozvoníl brátu} — When I had finished the book, I rang my brother.
\end{quote}

First the book was finished, then the call. \emph{\ru{Брату}} is \emph{\ru{брат}} with \textbf{-\ru{у}}: a man you ring takes that ending, and Mum would be \emph{\ru{маме}}.
\end{grammarlens}

\begin{grammarlens}[title={Grammar Lens: \ru{пока}, for two things at once}]
The goodbye \textbf{\ru{пока}} has a second job: \textquotedblleft{}while, for as long as\textquotedblright{}. Inside a story it usually takes the imperfective, and it lays two scenes side by side:

\begin{quote}
\textbf{\ru{Пока} \ru{я} \ru{читал}, \ru{брат} \ru{спал}.} — \emph{poká ya chitál, brat spal} — While I was reading, my brother was asleep.
\end{quote}

Neither half is finished; both went on through the same hour.
\end{grammarlens}

\begin{cousinweb}
\textbf{\ru{пока}} is \emph{\ru{по}-} with an old pronoun meaning \textquotedblleft{}which, what kind\textquotedblright{}, the ancestor of \emph{\ru{какой}}, on the question-root PIE \emph{*k\textsuperscript{w}o-} of \emph{\ru{как}} and of English \emph{who} and \emph{when}. By the usual account \textquotedblleft{}for as long as\textquotedblright{} is its old sense, and the goodbye, \textquotedblleft{}for now\textquotedblright{}, grew out of it in the twentieth century. \textbf{\ru{Когда}} comes off the same root.
\end{cousinweb}

\subsection*{Your turn}
Say these aloud:

\begin{itemize}
\raggedright
  \item \emph{Kogdá ya chitál, pozvoníl brat.} One call cut in
  \item \emph{Kogdá ya prochitál knígu, ya pozvoníl brátu.} One after the other
  \item \emph{Poká ya chitál, brat spal.} Two at once
  \item a woman's version: \emph{Poká ya chitála, brat spal.}
  \item one of your own: \emph{Poká ya …, …}, two things that went on together
\end{itemize}

\begin{tcolorbox}[breakable,colback=teal!4,colframe=teal!35!black,title={Before you move on}]
Which word puts two scenes side by side? (\emph{\textbf{\ru{Пока}.}}) \textquotedblleft{}When I was reading, my brother rang\textquotedblright{}: which verb is perfective? (\emph{\textbf{\ru{Позвонил}}}, the call that cut in.) Perfective on both sides of \emph{\ru{когда}}: at once, or one after the other? (\textbf{One after the other.})
\end{tcolorbox}
\section[(review)]{Review — what went on, and what got done}

\label{lesson:RU-C261-review}



\begin{quote}
Say \textquotedblleft{}I was at the museum\textquotedblright{}, as a man, then as a woman. (\emph{Ya byl v muzéye. Ya bylá v muzéye.}) Then \textquotedblleft{}we were at the museum\textquotedblright{}. (\emph{My býli v muzéye.})
\end{quote}

\begin{grammarlens}[title={What you've built}]
Three lessons, and Russian can now look back. The past swaps the infinitive's \emph{-\ru{ть}} for \emph{-\ru{л}}, agreeing with who did it. Nearly every verb in it shows one of two views: the \textbf{imperfective} for what went on, again and again or at all; the \textbf{perfective} for one deed done, with a result. \textbf{\ru{Пока}} lays two scenes side by side; \textbf{\ru{когда}} with a perfective marks the moment that cut in, or the next step.

One evening, told both ways:

\begin{quote}
\textbf{\ru{Пока} \ru{мама} \ru{ужинала}, \ru{папа} \ru{смотрел} \ru{фильм}.} — \emph{poká máma úzhinala, pápa smotrél film} — While Mum was having dinner, Dad was watching a film. \textbf{\ru{Когда} \ru{мама} \ru{поужинала}, \ru{она} \ru{позвонила} \ru{бабушке}.} — \emph{kogdá máma poúzhinala, oná pozvoníla bábushke} — When Mum had finished her dinner, she rang Grandma.
\end{quote}

The first pair went on together. The second pair came one after the other, each one finished: \emph{\ru{поужинала}} is the perfective of \emph{\ru{ужинать}}, made with \emph{\ru{по}-}.
\end{grammarlens}

\subsection*{Your turn}
\begin{itemize}
\raggedright
  \item \emph{Say it:} \emph{Poká máma úzhinala, pápa smotrél film.}
  \item \emph{Say it:} \emph{Kogdá ya chitál, pozvoníl brat.}
  \item \emph{Say it:} \emph{Kogdá my posmotréli film, my pozvoníli máme.}
  \item \emph{Answer:} \emph{Ty chitál étu knígu?} \emph{Da, chitál.}
  \item \emph{Say it:} your own evening: one thing that went on, with \emph{poká}, and one that got done
\end{itemize}

\begin{tcolorbox}[breakable,colback=teal!4,colframe=teal!35!black,title={Before you move on}]
\textquotedblleft{}She was reading\textquotedblright{}? (\emph{\textbf{\ru{Она} \ru{читала}.}}) \textquotedblleft{}We watched the film to the end\textquotedblright{}? (\emph{\textbf{\ru{Мы} \ru{посмотрели} \ru{фильм}.}}) Two scenes at once: which word? (\emph{\textbf{\ru{Пока}.}})
\end{tcolorbox}
